%============================================================
% PREÁMBULO OFICIAL FLUIDODINAMICA_BASE
%============================================================

\documentclass[12pt, a4paper]{article}

%------------------------------------------------------------
% Idioma, codificación y tipografía
%------------------------------------------------------------
\usepackage[spanish]{babel}
\usepackage[utf8]{inputenc}
\usepackage[T1]{fontenc}
\usepackage{lmodern}        % Tipografía moderna, clara y apta para PDF
\usepackage{microtype}      % Optimiza microtipografía (ideal para libros y apuntes)

%------------------------------------------------------------
% Matemática
%------------------------------------------------------------
\usepackage{amsmath, amssymb, amsfonts}
\usepackage{bm}             % Vectores y tensores en negrita: \bm{u}
\usepackage{mathtools}      % Extensiones de amsmath
\usepackage{siunitx}        % Notación correcta de unidades SI
\sisetup{
  locale = DE,            % usa coma decimal si querés; cambiar a US para punto
  per-mode = symbol,
  exponent-product = \cdot,
  separate-uncertainty = true
}

%------------------------------------------------------------
% Gráficos y figuras
%------------------------------------------------------------
\usepackage{graphicx}
\usepackage{float}
\usepackage{subcaption}     % Figuras con subfiguras
\usepackage{wrapfig}

%------------------------------------------------------------
% Diseño de página
%------------------------------------------------------------
\usepackage{geometry}
\geometry{
  a4paper,
  left=25mm,
  right=25mm,
  top=25mm,
  bottom=25mm
}

\usepackage{setspace}       % Control de espaciado (opcional)
\onehalfspacing             % Espaciado 1.5 (queda excelente para apuntes)

%------------------------------------------------------------
% Tablas profesionales
%------------------------------------------------------------
\usepackage{booktabs}       % Líneas horizontales profesionales
\usepackage{multirow}
\usepackage{tabularx}
\usepackage{longtable}

%------------------------------------------------------------
% Enlaces
%------------------------------------------------------------
\usepackage[hidelinks]{hyperref}  % Enlaces sin recuadros

%------------------------------------------------------------
% Colores y cajas
%------------------------------------------------------------
\usepackage{xcolor}
\usepackage{tcolorbox}
\tcbset{
  colback = white,
  colframe = black,
  arc = 2mm,
  boxrule = 0.6pt
}

%------------------------------------------------------------
% Ambientes personalizados
%------------------------------------------------------------

% Entorno para ejercicios
\newenvironment{ejercicio}[1][]{
  \begin{tcolorbox}[title={\textbf{Ejercicio #1}}, colback=gray!5]
}{
  \end{tcolorbox}
}

% Entorno para soluciones
\newenvironment{solucion}{
  \begin{tcolorbox}[title={\textbf{Solución}}, colback=blue!3]
}{
  \end{tcolorbox}
}

% Para dejar comentarios o notas al margen
\newcommand{\nota}[1]{\textcolor{blue}{\textbf{#1}}}

%============================================================
% FIN DEL PREÁMBULO
%============================================================
\begin{document}

%============================================================
% Ejercicio c.3 – Venturi + válvula de control en línea de proceso
%============================================================

\subsection*{Ejercicio c.3: Medición y control de caudal con Venturi y válvula}

Se requiere regular el caudal de agua que ingresa a un reactor que opera a
presión atmosférica. El sistema de control previsto incluye un medidor Venturi,
cuya señal de presión diferencial es enviada a un controlador que actúa sobre la
posición del obturador de una válvula de regulación. El caudal debe poder
ajustarse en el rango $10\text{–}50\ \text{m}^3/\text{h}$.

Las curvas características de tres válvulas candidatas (A, B y C) son provistas
por el fabricante y expresan el porcentaje del caudal máximo que pasa por la
válvula en función de la \% de apertura, \emph{para una caída de presión fija}
a través de la válvula. Las constantes de pérdida de carga para las válvulas
totalmente abiertas son $K_A=1$, $K_B=6$ y $K_C=7$.

\bigskip

\noindent\textbf{Datos adicionales:}
\begin{itemize}
  \item La presión en la sección del manómetro varía entre 2 y 3 bar(g).
  \item Longitud equivalente de la tubería: $L_e = 75\ \text{m}$.
  \item Tubería: diámetro nominal $3''$ Sch40, rugosidad $5\times 10^{-5}\ \text{m}$.
  \item Venturi: relación de diámetros $\beta = 0{,}4$, diámetro máximo igual al de la tubería.
  \item Propiedades del agua: $\rho = 1000\ \text{kg/m}^3$, $\mu = 1\times 10^{-3}\ \text{Pa\,s}$.
\end{itemize}

\noindent\textbf{Incógnitas:}
\begin{itemize}
  \item[(a)] Rango de lectura del manómetro diferencial del Venturi para $10\text{–}50\ \text{m}^3/\text{h}$.
  \item[(a)] Calidad cualitativa de la señal para un lazo de control de caudal.
  \item[(b)] Rango de apertura de las válvulas A, B y C para cubrir el rango de operación.
  \item[(b)] Selección o descarte de alguna válvula en base a su característica.
\end{itemize}

%============================================================
\subsubsection*{1. Hipótesis de modelado}
\begin{enumerate}
  \item Flujo estacionario, agua incompresible.
  \item Flujo turbulento en la línea (se verificará con Reynolds).
  \item Venturi operado en condiciones de líquido → factor de expansibilidad $Y=1$.
  \item Coeficiente de descarga del Venturi según teórico: 
        \[
        C_D = 0{,}9858 - 0{,}196\,\beta^{4{,}5}.
        \]
  \item Válvula modelada mediante pérdida de carga cuadrática:
        \[
        \Delta P = C\,K_v\,Q^2,
        \]
        válida para válvulas en servicio turbulento.
  \item Curvas del fabricante: $\Delta P = \text{cte}$ para todos los puntos → 
        se cumple
        \[
        \frac{Q_x}{Q_{\max}}
        = \sqrt{\frac{K_{v,\mathrm{TA}}}{K_{v,x}}},
        \]
        lo que justifica el uso de \%Q\(_{\max}\) como ratio de caudal relativo.
\end{enumerate}

%============================================================
\subsubsection*{2. Geometría del Venturi y áreas}
La tubería $3''$ Sch40 tiene diámetro interno
\[
D \approx 0{,}078\ \text{m},
\qquad
A_1 = \frac{\pi D^2}{4}
    \approx 4{,}77\times 10^{-3}\ \text{m}^2.
\]

El estrechamiento del Venturi es ($\beta = 0{,}4$)
\[
d = \beta D = 0{,}4\cdot 0{,}078 \approx 0{,}031\ \text{m},
\qquad
A_2 = \frac{\pi d^2}{4} \approx 7{,}65\times 10^{-4}\ \text{m}^2,
\]
\[
1-\beta^4 = 1 - (0{,}4)^4 = 0{,}9744,
\qquad
C_D = 0{,}9858 - 0{,}196\beta^{4{,}5} \approx 0{,}9826.
\]

%============================================================
\subsubsection*{3. Parte (a): Lectura del manómetro del Venturi}

La ecuación reducida de Venturi para líquido incompresible es
\[
Q = C_D\,A_2
\sqrt{\frac{2\,\Delta P}{\rho\,(1-\beta^4)}},
\]
de donde
\[
\Delta P
= 
\frac{\rho (1-\beta^4)}{2}
\left(\frac{Q}{C_D A_2}\right)^2.
\]

Convertimos el rango de caudales del problema:
\[
Q_{\min} = \frac{10}{3600} = 2{,}78\times 10^{-3}\ \text{m}^3/\text{s},
\]
\[
Q_{\max} = \frac{50}{3600} = 1{,}39\times 10^{-2}\ \text{m}^3/\text{s}.
\]

\paragraph{Presión diferencial mínima.}
\[
\Delta P_{\min} =
\frac{1000\cdot 0{,}9744}{2}
\left(\frac{2{,}78\times 10^{-3}}
           {0{,}9826\cdot 7{,}65\times 10^{-4}}\right)^2
\approx 6{,}66\times 10^{3}\ \text{Pa}.
\]

\paragraph{Presión diferencial máxima.}
\[
\Delta P_{\max}
= \Delta P_{\min}
  \left(\frac{Q_{\max}}{Q_{\min}}\right)^2
= 6{,}66\times 10^3 \left(\frac{50}{10}\right)^2
\approx 1{,}67\times 10^{5}\ \text{Pa}.
\]

\[
\boxed{
6{,}7\ \text{kPa}\ \lesssim\ \Delta P\ \lesssim\ 167\ \text{kPa}
}
\]


\paragraph{Comentario sobre la calidad de la regulación.}

Dado que 
\[
\Delta P \propto Q^2,
\]
la sensibilidad en caudal es muy pobre en el extremo bajo del rango:
mientras que $Q$ cambia por un factor de 5, la $\Delta P$ cambia por un factor 25.

De aquí que:

- Para $Q\approx 10\ \text{m}^3/h$, el transmisor opera muy cerca de cero → mala relación señal/ruido.
- Para $Q\approx 50\ \text{m}^3/h$, la señal está cerca del 80 \% de escala → excelente resolución.

De no usarse un transmisor con extracción de raíz cuadrada, la respuesta del lazo
sería extremadamente no lineal.

\[
\text{Conclusión: el Venturi tiene buena sensibilidad media–alta, pobre en caudales bajos.}
\]

%============================================================
\subsubsection*{4. Parte (b): Válvulas de control y rango de apertura}

Tomamos la ecuación general de pérdida de carga para válvulas operando
en régimen turbulento, modelo Crane:

\[
\Delta P = C\,K_v\,Q^2.
\]

Dado que las curvas del fabricante están construidas a $\Delta P=\text{cte}$, para dos
condiciones de apertura $x$ y totalmente abierta (TA):

\[
C\,K_{v,x}\,Q_x^2 = C\,K_{v,\mathrm{TA}}\,Q_{\max}^2.
\]

entonces:

\[
\frac{Q_x}{Q_{\max}}
= 
\sqrt{\frac{K_{v,\mathrm{TA}}}{K_{v,x}}}.
\]

Esto justifica la definición de eje vertical en las curvas del fabricante:
\[
\%\ Q_{\max} = \frac{Q_x}{Q_{\max}}\times 100.
\]

\bigskip

\paragraph{Interpretación de las curvas características.}

A partir del gráfico:

\begin{itemize}
  \item \textbf{Válvula A} (quick-opening):  
        El $20$\% de $Q_{\max}$ se alcanza con aperturas del orden $5$–$8$\%.  
        Es extremadamente sensible a bajas aperturas.  
        No apta para control modulante. Se descarta.
  \item \textbf{Válvula B} (característica lineal):  
        El $20$\% de caudal se obtiene en torno al $20$\% de apertura.  
        Adecuada para control razonablemente lineal.  
        Rango operativo útil: $20$–$100$\%.
  \item \textbf{Válvula C} (equal-percentage):  
        El $20$\% de $Q_{\max}$ se obtiene cerca de $40$–$45$\% de apertura.  
        Permite controlar un rango amplio gracias a su comportamiento exponencial.  
        Es la preferida en control de caudal con rangeabilidad 5:1.
\end{itemize}

\paragraph{Conclusión de selección.}

\[
\boxed{
\text{A se descarta; B es aceptable; C es la más adecuada para control.}
}
\]

La válvula C ofrece:

- mejor distribución de sensibilidad sobre el rango 10–50 m$^3$/h,  
- menor saturación en aperturas bajas,  
- comportamiento estable ante variaciones de $\Delta P$,  
- y mayor “rangeabilidad instalada”, como indican los apuntes de clase de control.

\end{document}

